% Chapter introduction: why the topic arose, who first asked, what it answers.
% STATUS: written 2026-11-04 (sources of the chapter read in full, each from the page)
Every associative ring carries a second product next to its own: the commutator
$[a,b]=ab-ba$. A map that obeys the product rule for the commutator is a Lie
derivation, and Martindale began the theory in 1964 with a single question
\cite[p.~183]{Martindale1964}. Every derivation is a Lie derivation, and so is
every sum of a derivation and a central map that sends commutators to zero. The
question is whether there is anything else.

This chapter follows the answers class by class. Martindale gave the first one
for primitive rings with a nontrivial idempotent
(Section~\ref{sec:lie-type:martindale}). Brešar settled the question for
noncommutative prime rings through his theory of functional identities
(Section~\ref{sec:lie-type:bresar}). Iterating the bracket gives Lie triple and
$n$-Lie derivations (Section~\ref{sec:lie-type:triple-and-n}). Dropping
linearity gives the nonlinear maps of
Section~\ref{sec:lie-type:nonlinear}. Section~\ref{sec:lie-type:lie-on-algebras}
collects the answers on the algebras where the question can be asked
in explicit terms, from $B(X)$ to matrix and nest algebras.

\hbfigure{05-lie-type/standard-form}{The question of Martindale: is every Lie derivation the sum of a derivation $D$ and a central map $\tau$ that vanishes on commutators?}{fig:lie:standard}
